\documentclass[11pt]{amsart}
\usepackage{amsmath,amssymb,amsthm,mathtools}
\usepackage{geometry}
\usepackage{booktabs}
\usepackage{xcolor}
\usepackage{hyperref}
\usepackage{enumitem}
\geometry{margin=1in}
\hypersetup{colorlinks=true,linkcolor=blue!70!black,citecolor=blue!70!black,urlcolor=blue!70!black,
  pdftitle={Hodge-index and hyperbolicity reformulations of the Riemann hypothesis as positivity statements},
  pdfauthor={David Alejandro Trejo Pizzo}}

\newtheorem{theorem}{Theorem}[section]
\newtheorem{proposition}[theorem]{Proposition}
\newtheorem{lemma}[theorem]{Lemma}
\newtheorem{corollary}[theorem]{Corollary}
\newtheorem{conjecture}[theorem]{Conjecture}
\theoremstyle{definition}
\newtheorem{definition}[theorem]{Definition}
\theoremstyle{remark}
\newtheorem{remark}[theorem]{Remark}
\newtheorem{observation}[theorem]{Observation}
\newtheorem*{thesis}{Heuristic summary}

\newcommand{\re}{\operatorname{Re}}
\newcommand{\im}{\operatorname{Im}}
\newcommand{\half}{\tfrac12}
\newcommand{\bmin}{\beta_{\min}}
\newcommand{\LP}{\mathcal{LP}}

\title[Hodge-index and hyperbolicity reformulations]{Hodge-index and hyperbolicity reformulations\\[3pt]
of the Riemann hypothesis as positivity statements}
\author{David Alejandro Trejo Pizzo}
\thanks{Independent Researcher. \texttt{dtrejopizzo@gmail.com}}

\begin{document}
\begin{abstract}
We examine two reformulations of the Riemann hypothesis (RH) whose mechanism is not, at first sight, a
quadratic-form positivity: a Hodge-index picture built on the band-limited Weil--Carleson form, and
hyperbolicity of the Jensen polynomials of $\Xi(t)=\xi(\half+it)$. We record the moment identity
$b(k)=(-1)^ka(k)=2^{2k+3}m_{2k}/(2k)!$, expressing the Taylor coefficients of $\Xi$ through the even moments of the
de Bruijn--Newman kernel $\Phi$; we state the classical Hermite and Pólya-frequency criteria in this setting and
the Laguerre closure of the Laguerre--Pólya class under differentiation; and we give numerical observations on a
sine-kernel Gram matrix of zeta zeros, together with an elementary upper bound for its smallest eigenvalue in
terms of the smallest normalized gap. We argue heuristically that both routes reduce to positivity statements.
\end{abstract}

\maketitle

\tableofcontents
\newpage

\section{Introduction}\label{sec:intro}
In \cite{P12} five analytic reformulations of RH were compared, and \cite{P13} treats the de Bruijn--Newman
zero flow. We take from them the following heuristic
theme: the unconditional tools available for the nontrivial zeros are positivity statements, which tend
to yield lower bounds or existence results, whereas RH is a one-sided constraint (no zeros off the line).

This paper looks at two structures that might seem to escape that theme. The first is the
Hodge-index mechanism, by which the Riemann hypothesis for curves over finite fields
\cite{Weil,MattuckTate,Groth} obtains an upper bound on Frobenius eigenvalues from a positivity via duality. The second is
hyperbolicity (real-rootedness), studied in combinatorics through interlacing families
\cite{MSS1,MSS} and stability preservers \cite{BB1,BB2}. In both cases the relevant criteria are classical
positivity statements (Hermite's criterion, total positivity), and we record them precisely in the present
setting.

What is proved here: the moment identity \eqref{eq:moment}; the elementary bound of
Lemma~\ref{thm:reggap}; Proposition~\ref{prop:blind} (a consequence of Laguerre's theorem). The criteria
of \S\ref{sec:hyp} are classical and cited. The remaining statements are numerical observations, reproducible
with the script \texttt{gram.py}, or heuristic remarks, labelled as such.

\section{The Hodge-index picture}\label{sec:hodge}

\subsection{The band-limited form}
Fix a height $T_0$, a band $d$ and an $N$-dimensional sinc span of test functions with Fourier support in
$[-d,d]$, as in \cite{P12}. Let $\mathfrak t=A-P$, where $A$ and $P$ are the Gram matrices on that span of the
archimedean density and of the prime term of the Weil explicit formula (denoted $A_\Phi$, $P_F$ in \cite{P12}).
Numerically, at $T_0=1000$ and $N=29$, the smallest eigenvalue of $\mathfrak t$ relative to $A$, namely
$1-\lambda_{\max}(P,A)$, decreases from $0.66$ at $d=1$ to $0$ within five digits at $d=3$ and $d=5$
\cite[Observation ``Numerical saturation'']{P12}. We view
this heuristically as an analogue of a degenerating intersection form: the Hodge index theorem for a surface
$C\times C$ uses a definite primitive form, while here the candidate form appears to lose its gap. No
implication with the de Bruijn--Newman statement $\Lambda=0$ is claimed.

\subsection{A sine-kernel Gram matrix of the window zeros}
Heuristically, the null directions of $\mathfrak t$ are band-limited functions vanishing at the zeros in the
window, and the natural regularized quantity is the Gram matrix of point evaluations at those zeros in the
Paley--Wiener space. We do not prove an identification of the smallest positive eigenvalue of $\mathfrak t$
with the quantity below; we study the Gram matrix directly.

Let $\gamma_1<\dots<\gamma_K$ be $K$ consecutive ordinates of zeros of $\zeta$ above $T_0$, normalized to mean
spacing one by $x_i=\gamma_i\log(T_0/2\pi)/(2\pi)$, and let $\bmin=\min_i(x_{i+1}-x_i)$.

\begin{lemma}\label{thm:reggap}
Let $G_{ij}=\dfrac{\sin\pi(x_i-x_j)}{\pi(x_i-x_j)}$ $(G_{ii}=1)$. Then $G\succ0$, and
\begin{equation}\label{eq:gapident}
\lambda_{\min}(G)\;\le\;1-\frac{\sin\pi\bmin}{\pi\bmin}\;\le\;\frac{\pi^2}{6}\,\bmin^2 .
\end{equation}
\end{lemma}

\begin{proof}
$G_{ij}=\int_{-1/2}^{1/2}e^{2\pi i(x_i-x_j)\xi}\,d\xi$, so $c^*Gc=\int_{-1/2}^{1/2}|\sum_ic_ie^{-2\pi ix_i\xi}|^2d\xi$,
which vanishes only if the exponential sum vanishes identically, impossible for distinct $x_i$ and $c\ne0$.
If $x_j-x_i=\bmin$, the $2\times2$ principal submatrix on $\{i,j\}$ has eigenvalues $1\pm\operatorname{sinc}$,
with $\operatorname{sinc}=\sin(\pi\bmin)/(\pi\bmin)$; by Cauchy interlacing $\lambda_{\min}(G)\le1-\operatorname{sinc}$.
Finally $\sin y\ge y-y^3/6$ for $y\ge0$ gives $1-\sin y/y\le y^2/6$ with $y=\pi\bmin$.
\end{proof}

Positivity of $G$ holds for any distinct points and has no bearing on RH. Numerically (Table~\ref{tab:gram},
$K=25$, ordinates from \texttt{mpmath.zetazero}, script \texttt{gram.py}) the ratio
$\lambda_{\min}(G)/\bmin^2$ lies in $[0.06,0.63]$, well below the bound $\pi^2/6\approx1.645$ of
\eqref{eq:gapident}, and varies with $T_0$: for $K=25$ the smallest eigenvalue is governed by the whole
configuration, not only by the closest pair, and no asymptotic relation $\lambda_{\min}(G)\asymp\bmin^2$ is
supported by these data.

\begin{table}[ht]
\centering
\begin{tabular}{rrrrr}
\toprule
$T_0$ & $\lambda_{\min}(G)$ & $\bmin$ & $1-\frac{\sin\pi\bmin}{\pi\bmin}$ & $\lambda_{\min}(G)/\bmin^2$ \\
\midrule
$100$ & $0.0094$ & $0.3722$ & $0.2128$ & $0.068$ \\
$1000$ & $0.0672$ & $0.3568$ & $0.1967$ & $0.528$ \\
$3000$ & $0.0593$ & $0.3452$ & $0.1849$ & $0.497$ \\
$10^4$ & $0.0609$ & $0.3125$ & $0.1531$ & $0.624$ \\
\bottomrule
\end{tabular}
\caption{Sine-kernel Gram of $K=25$ consecutive zeros above $T_0$ (output of \texttt{gram.py}).}\label{tab:gram}
\end{table}

\begin{remark}\label{cor:price}
By \eqref{eq:gapident}, a uniform lower bound $\lambda_{\min}(G)\ge c>0$ (over windows) implies uniform
separation $\bmin\ge(6c)^{1/2}/\pi$ of the normalized zeros. No implication between such a bound and RH is
known in either direction (the bound concerns the spacing of the ordinates, not the real parts), and the bound
is expected to fail because of the small gaps predicted by GUE statistics.
\end{remark}

\begin{remark}[heuristic]\label{rem:hodgecap}
In frame language, for a fixed number $K$ of points the upper bound $\lambda_{\max}(G)\le K$ is trivial, while a
lower frame bound means that no two zeros are too close, a one-sided constraint on clustering. Selberg's
unconditional results on $S(T)=\frac1\pi\arg\zeta(\half+iT)$ concern its typical size
($\sqrt{\log\log T}$), not its extremal behaviour. For Riesz sequences of reproducing kernels see \cite{Pavlov,Seip}.
\end{remark}

\section{The hyperbolicity route}\label{sec:hyp}

\subsection{The criterion and the moment identity}
Let $\Xi(t)=\xi(\half+it)$, where $\xi(s)=\tfrac12s(s-1)\pi^{-s/2}\Gamma(s/2)\zeta(s)$, and write
\[\Xi(t)=\sum_{k\ge0}a(k)t^{2k},\qquad a(k)=\frac{\Xi^{(2k)}(0)}{(2k)!}.\]
In the variable $z=-t^2$
\[
\Xi=\sum_{k\ge0}\frac{\gamma(k)}{k!}z^k,\qquad \gamma(k)=k!\,b(k),\quad b(k)=(-1)^ka(k).
\]
The Laguerre--P\'olya class $\LP$ consists of the locally uniform limits of real polynomials with only real
zeros; RH holds if and only if $\Xi\in\LP$, equivalently $F(z)=\sum_k\gamma(k)z^k/k!\in\LP$. The Jensen
polynomials are
\begin{equation}\label{eq:jensen}
J^{d,n}(X)=\sum_{j=0}^d\binom{d}{j}\gamma(n+j)X^j .
\end{equation}
RH holds if and only if $J^{d,n}$ is hyperbolic for all $d,n\ge0$ (Pólya \cite{Polya}; see \cite{GORZ}, who
prove hyperbolicity for each fixed $d$ and all sufficiently large $n$).

Let
\[
\Phi(u)=\sum_{n\ge1}\big(2\pi^2n^4e^{9u}-3\pi n^2e^{5u}\big)e^{-\pi n^2e^{4u}},
\]
an even function of $u$. With the normalization $H_0(z)=\int_0^\infty\Phi(u)\cos(zu)\,du=\frac18\Xi(z/2)$
of de Bruijn and Newman (as in Rodgers--Tao \cite{RT}) one has $\Xi(t)=8\int_0^\infty\Phi(u)\cos(2tu)\,du$.
Differentiating $2k$ times under the integral (justified by the super-exponential decay of $\Phi$) gives
\begin{equation}\label{eq:moment}
b(k)=(-1)^k a(k)=\frac{2^{2k+3}\,m_{2k}}{(2k)!},\qquad m_{2k}=\int_0^\infty u^{2k}\Phi(u)\,du .
\end{equation}
Thus the Taylor coefficients of $\Xi$ are, up to sign and the factor $2^{2k+3}/(2k)!$, the even moments of $\Phi$.
By the Aissen--Schoenberg--Whitney--Edrei theorem \cite{Edrei} (see also \cite{Karlin}), RH holds if and only
if $\{b(k)\}$ is a Pólya frequency sequence, i.e.\ the Toeplitz matrix $[b(i-j)]_{i,j\ge0}$ (with $b(k)=0$ for
$k<0$) is totally positive, i.e.\ all its minors are nonnegative.

Numerically (\texttt{gram.py}, $0\le k\le8$) the identity \eqref{eq:moment} agrees with the Taylor coefficients
of $\Xi$ for $k\le4$, and the inequalities $b(k)^2>b(k-1)b(k+1)$, $1\le k\le7$, hold; they follow from the
Turán inequalities $\gamma(k)^2\ge\gamma(k-1)\gamma(k+1)$, proved by Csordas, Norfolk and Varga \cite{CNV}.

\subsection{Differentiation preserves the Laguerre--Pólya class}
In terms of $\Phi$, the map $\Phi\mapsto4u^2\Phi$ corresponds to $\Xi\mapsto-\Xi''$. In the variable $z$, write
$F(z)=\sum_k\gamma(k)z^k/k!$, so that $\Xi(t)=F(-t^2)$; the shift $n\mapsto n+1$ of the Jensen polynomials
\eqref{eq:jensen} replaces $\gamma(k)$ by $\gamma(k+1)$, i.e.\ $F$ by $F'$.

\begin{proposition}\label{prop:blind}
If $f\in\LP$ then $-f''\in\LP$. Likewise, if $F\in\LP$ then $F'\in\LP$; hence the shift $n\mapsto n+1$ maps the
Jensen polynomials of any $\LP$ function to those of another $\LP$ function, and it preserves hyperbolicity
independently of any arithmetic input.
\end{proposition}

\begin{proof}
By Laguerre's theorem, $\LP$ is closed under differentiation \cite[Ch.~VIII]{Levin}; it is also closed under
multiplication by $-1$. This gives both closure statements. The Jensen polynomials of $F'$ are
$\sum_j\binom dj\gamma(n+1+j)X^j=J^{d,n+1}$, and by Pólya's theorem \cite{Polya} all Jensen polynomials of an
$\LP$ function are hyperbolic.
\end{proof}

\begin{remark}[heuristic]
Since the same closure holds for every $\LP$ function (for instance $\cos t$), operators of this kind do not
distinguish $\Xi$; any arithmetic input must enter through the statement $\Xi\in\LP$ itself, i.e.\ through the
moments $\{m_{2k}\}$.
\end{remark}

\subsection{Real-rootedness as a positivity}
\begin{theorem}[Hermite]\label{thm:hermite}
A real polynomial of degree $d$ is real-rooted if and only if its Hankel matrix $(p_{i+j})_{i,j=0}^{d-1}$ of
power sums of the roots is positive semidefinite \cite{Hermite,Obreschkoff}.
\end{theorem}

\begin{remark}[heuristic]
Hence real-rootedness of each $J^{d,n}$ is a quadratic positivity. The interlacing method of
\cite{MSS1,MSS} uses real stability of $\det(xI+\sum_iz_iA_i)$ for $A_i\succeq0$, proved in \cite{BB2},
again a statement built from positive semidefinite matrices. Numerically (\texttt{gram.py}) the Hermite matrices of
$J^{d,n}$ for $2\le d\le4$, $0\le n\le2$ are positive definite, in agreement with the theorem of
\cite{GORZ} that $J^{d,n}$ is hyperbolic for $d\le8$ and all $n$.
\end{remark}

\section{Summary}\label{sec:capstone}

\begin{thesis}
In the two settings above the available criteria are positivity statements: Hermite's criterion and total
positivity for hyperbolicity, and Gram-matrix bounds for the Hodge-index picture, where a uniform lower bound on
$\lambda_{\min}(G)$ implies separation of the zeros (Remark~\ref{cor:price}) and is not known to be related to RH. We regard this as
heuristic support for the view that positivity-based methods yield one-sided information of the wrong kind
for RH; no theorem about methods is claimed.
\end{thesis}

\section{Conclusion}\label{sec:concl}
We recorded the moment identity \eqref{eq:moment}, the classical Hermite and Pólya-frequency criteria in this
setting, Laguerre's closure (Proposition~\ref{prop:blind}), the elementary bound of Lemma~\ref{thm:reggap}, and numerical observations on a
sine-kernel Gram (Table~\ref{tab:gram}); we argued heuristically that both routes reduce to positivity statements.

\begin{thebibliography}{99}
\bibitem{P12} D.~A.~Trejo Pizzo, \emph{The Band-Limited Weil--Carleson Form: a numerically saturated generalized eigenvalue, sine-kernel structure of its spectrum, and a comparison of five reformulations}, Paper 12 of this series.
\bibitem{P13} D.~A.~Trejo Pizzo, \emph{Monotonicity of the Off-Line Mass under the de Bruijn--Newman Zero Flow}, Paper 13 of this series.
\bibitem{BB1} J.~Borcea, P.~Brändén, \emph{The Lee--Yang and Pólya--Schur programs. II. Theory of stable polynomials and applications}, Comm.\ Pure Appl.\ Math.\ \textbf{62} (2009), 1595--1631.
\bibitem{BB2} J.~Borcea, P.~Brändén, \emph{Applications of stable polynomials to mixed determinants: Johnson's conjectures, unimodality, and symmetrized Fischer products}, Duke Math.\ J.\ \textbf{143} (2008), 205--223.
\bibitem{CNV} G.~Csordas, T.~S.~Norfolk, R.~S.~Varga, \emph{The Riemann hypothesis and the Turán inequalities}, Trans.\ Amer.\ Math.\ Soc.\ \textbf{296} (1986), 521--541.
\bibitem{Edrei} A.~Edrei, \emph{On the generating functions of totally positive sequences. II}, J.\ Analyse Math.\ \textbf{2} (1952), 104--109; see also A.~Edrei, \emph{Proof of a conjecture of Schoenberg on the generating function of a totally positive sequence}, Canad.\ J.\ Math.\ \textbf{5} (1953), 86--94.
\bibitem{Groth} A.~Grothendieck, \emph{Sur une note de Mattuck--Tate}, J.\ Reine Angew.\ Math.\ \textbf{200} (1958), 208--215.
\bibitem{GORZ} M.~Griffin, K.~Ono, L.~Rolen, D.~Zagier, \emph{Jensen polynomials for the Riemann zeta function and other sequences}, Proc.\ Natl.\ Acad.\ Sci.\ USA \textbf{116} (2019), 11103--11110.
\bibitem{Hermite} C.~Hermite, \emph{Sur le nombre des racines d'une équation algébrique comprises entre des limites données}, J.\ Reine Angew.\ Math.\ \textbf{52} (1856), 39--51.
\bibitem{Pavlov} S.~V.~Hru\v{s}\v{c}ev, N.~K.~Nikolskii, B.~S.~Pavlov, \emph{Unconditional bases of exponentials and of reproducing kernels}, Lecture Notes in Math.\ \textbf{864} (1981), 214--335.
\bibitem{Karlin} S.~Karlin, \emph{Total Positivity}, Vol.~I, Stanford Univ.\ Press, 1968.
\bibitem{Levin} B.~Ja.~Levin, \emph{Distribution of Zeros of Entire Functions}, Transl.\ Math.\ Monogr.\ \textbf{5}, AMS, 1964.
\bibitem{MattuckTate} A.~Mattuck, J.~Tate, \emph{On the inequality of Castelnuovo--Severi}, Abh.\ Math.\ Sem.\ Univ.\ Hamburg \textbf{22} (1958), 295--299.
\bibitem{MSS1} A.~Marcus, D.~Spielman, N.~Srivastava, \emph{Interlacing families I: Bipartite Ramanujan graphs of all degrees}, Ann.\ of Math.\ \textbf{182} (2015), 307--325.
\bibitem{MSS} A.~Marcus, D.~Spielman, N.~Srivastava, \emph{Interlacing families II: Mixed characteristic polynomials and the Kadison--Singer problem}, Ann.\ of Math.\ \textbf{182} (2015), 327--350.
\bibitem{Obreschkoff} N.~Obreschkoff, \emph{Verteilung und Berechnung der Nullstellen reeller Polynome}, VEB Deutscher Verlag der Wissenschaften, Berlin, 1963.
\bibitem{Polya} G.~Pólya, \emph{Über die algebraisch-funktionentheoretischen Untersuchungen von J.~L.~W.~V.~Jensen}, Kgl.\ Danske Vid.\ Sel.\ Math.-Fys.\ Medd.\ \textbf{7} (1927), 3--33.
\bibitem{RT} B.~Rodgers, T.~Tao, \emph{The de Bruijn--Newman constant is non-negative}, Forum Math.\ Pi \textbf{8} (2020), e6.
\bibitem{Seip} K.~Seip, \emph{Interpolation and Sampling in Spaces of Analytic Functions}, Univ.\ Lecture Ser.\ \textbf{33}, AMS, 2004.
\bibitem{Weil} A.~Weil, \emph{Sur les courbes algébriques et les variétés qui s'en déduisent}, Hermann, Paris, 1948.
\end{thebibliography}

\end{document}
